\section{Visibility Graphs}

A visibility graph converts a continuous geometric planning problem into a
discrete graph-search problem.

For polygonal obstacles, the graph will eventually contain the start point,
goal point, and obstacle vertices as nodes. Two nodes are connected when they
are geometrically visible to one another.

\subsection{Geometric Visibility}

Two points $A$ and $B$ are considered visible when the open segment between
them does not pass through the interior of any obstacle.

For an obstacle $O$, this can be expressed as

\[
(A,B) \cap \operatorname{Interior}(O) = \varnothing.
\]

The use of the open segment $(A,B)$ is important. A visibility segment may
terminate exactly at an obstacle vertex without entering the obstacle
interior.

Visibility therefore differs from the Phase 1 collision convention. In
ordinary collision checking, touching an obstacle boundary counts as a
collision. In visibility testing, touching or traveling along an obstacle
boundary may be allowed as long as the segment does not enter the obstacle
interior.

\subsection{Proper Boundary Crossings}

A candidate visibility segment is invalid if it properly intersects an
obstacle edge. A proper intersection crosses through the interiors of both
segments rather than merely touching at an endpoint.

Proper-intersection testing detects ordinary crossings through an obstacle,
but it is not sufficient by itself. A candidate segment can enter an obstacle
through one of its vertices without producing a proper intersection.

\subsection{Parametric Representation of a Segment}

A point along the segment from $A$ to $B$ can be represented parametrically
as

\[
P(t) = A + t(B-A),
\qquad 0 \leq t \leq 1.
\]

The parameter $t$ describes the position along the segment:

\[
P(0)=A,
\]

\[
P(1)=B,
\]

and

\[
P\left(\frac{1}{2}\right)
=
\frac{A+B}{2}.
\]

If a point $V$ is known to lie on the nonzero segment $\overline{AB}$, its
parameter can be recovered from either coordinate. For example,

\[
t = \frac{V_x-A_x}{B_x-A_x}
\]

when the $x$ component is suitable for division.

The implementation uses the coordinate with the larger absolute change
between $A$ and $B$. This avoids division by zero for vertical or horizontal
segments and reduces division by very small direction components.

The helper that computes this parameter assumes that the point already lies
on a nonzero segment. It does not perform segment-membership validation.

\subsection{Visibility Near Obstacle Vertices}

Proper-intersection testing alone cannot detect every invalid visibility
segment. In particular, a segment may pass through an obstacle vertex and
immediately enter the obstacle interior.

Suppose an obstacle vertex $V$ lies on a candidate segment and has parameter
$t_V$. Points immediately before and after the vertex can be sampled using

\[
P(t_V-\delta)
\]

and

\[
P(t_V+\delta),
\]

where $\delta$ is a small positive parameter displacement.

If either valid neighboring sample lies strictly inside the obstacle, then the
candidate segment enters the obstacle interior and is therefore not visible.

The endpoint cases require special treatment. If

\[
t_V=0,
\]

then $V=A$, so no portion of the candidate segment exists before the vertex.
Only the sample after the vertex needs to be considered. Similarly, if

\[
t_V=1,
\]

then $V=B$, so only the sample before the vertex needs to be considered.

The sampled parameters are also clamped to the valid segment interval

\[
0 \leq t \leq 1.
\]

This ensures that visibility testing considers only the finite candidate
segment rather than points on the infinite line extending beyond its
endpoints.

\subsection{Visibility Graph Construction}

Given a start point $S$, a goal point $G$, and a collection of polygonal
obstacles, the visibility graph is represented as

\[
\mathcal{G} = (V,E).
\]

The node set contains the start point, goal point, and every obstacle vertex:

\[
V =
\{S,G\}
\cup
\bigcup_{O \in \mathcal{O}}
\operatorname{Vertices}(O).
\]

In the current implementation, the start point is stored as node $0$, the
goal point as node $1$, and obstacle vertices are appended afterward.

To construct the edge set, every unique pair of nodes $(v_i,v_j)$ with

\[
i < j
\]

is tested for geometric visibility. If the two nodes are visible, an
undirected edge is added between them.

Thus,

\[
E =
\left\{
(v_i,v_j)
\mid
i < j
\text{ and }
\operatorname{Visible}(v_i,v_j)
\right\}.
\]

Using $i<j$ ensures that each unordered pair is considered exactly once. It
also prevents self-edges such as $(v_i,v_i)$.

\subsection{Euclidean Edge Weights}

Each visibility edge is assigned a weight equal to the Euclidean distance
between its endpoints.

For points

\[
A=(x_A,y_A)
\qquad\text{and}\qquad
B=(x_B,y_B),
\]

the edge weight is

\[
w(A,B)
=
\sqrt{
(x_B-x_A)^2
+
(y_B-y_A)^2
}.
\]

These weights represent the geometric length of traveling directly between
two visible nodes.

The resulting weighted graph can later be searched using shortest-path
algorithms such as Dijkstra's algorithm or A*.

\section{Shortest Paths with Dijkstra's Algorithm}

Once the weighted visibility graph has been constructed, motion planning
becomes a shortest-path problem on a graph.

For each node $v$, Dijkstra's algorithm maintains a tentative distance

\[
d(v),
\]

representing the shortest currently known distance from the start node to
$v$.

Initially,

\[
d(S)=0,
\]

while every other node is assigned

\[
d(v)=\infty.
\]

The central operation in Dijkstra's algorithm is \emph{edge relaxation}.
For an edge from $u$ to $v$ with weight $w(u,v)$, a new candidate distance
to $v$ is

\[
d(u)+w(u,v).
\]

If

\[
d(u)+w(u,v)<d(v),
\]

then a shorter route to $v$ has been discovered, and we update

\[
d(v)=d(u)+w(u,v).
\]

The predecessor of $v$ is also recorded as $u$ so that the final shortest
path can later be reconstructed.

\subsection{Selecting the Next Node}

After relaxing the edges of the current node, Dijkstra's algorithm selects
the unvisited node with the smallest tentative distance.

Once a node is selected in this way, its shortest-path distance is considered
finalized.

This property relies on all graph edge weights being non-negative. In the
visibility graph, edge weights are Euclidean distances, so this condition is
automatically satisfied.

\subsection{Path Reconstruction}

Whenever relaxation improves the tentative distance to a node $v$, the
algorithm records the node $u$ from which that improvement came.

This node is called the \emph{predecessor} or \emph{parent} of $v$.

After the goal is reached, the shortest path is reconstructed by following
the parent links backward:

\[
G \rightarrow \operatorname{parent}(G)
\rightarrow \operatorname{parent}(\operatorname{parent}(G))
\rightarrow \cdots \rightarrow S.
\]

Because this produces the path in reverse order, the resulting sequence is
reversed to obtain the final path from start to goal.

\subsection{Undirected Visibility Edges}

The visibility graph is undirected: if node $u$ can see node $v$, then node
$v$ can also see node $u$.

For efficiency, each undirected edge is stored only once in the graph.

Therefore, when processing a node, Dijkstra's algorithm must check whether
the current node appears as either endpoint of an edge.